% ARCHIVED DRAFT: root independently audited and incorporated 2026-11-19 17:55 UTC.
% Main part_B.tex is authoritative; this file preserves drafting provenance.
\begin{saresult}{B.079}{Proved}{Coefficient base change and an explicit monic saturation obstruction}
For an arbitrary unital inclusion $R\subset S$, let
$C=\bigcup_d(R:_S D_d(R))$ be the effective coefficient ring of B.064,
and put $Q=C/R$. Define the $R$-module directly by finite coefficient
relations
\[
 N(R,Q)=\left\{(q_0,q_1,\ldots)\in\bigoplus_{i\ge0}Q:
             \sum_i r^i q_i=0\text{ for every }r\in R\right\}.
                                                        \tag{B79.1}
\]
If the polynomial tensor map $\theta_2$ is injective, then $R\to C$
is epic and
\[
 C\otimes_R\Omega(R,S)\ \xrightarrow{\ \sim\ } C[X],
 \qquad C\otimes_RN(R,Q)=0.                              \tag{B79.2}
\]
The two assertions in (B79.2) are equivalent, without assuming
injectivity of $\theta_2$. They are necessary conditions, not a
universal sufficient criterion. Even when $C$ is a flat localization
of a domain $R$, they do not imply injectivity of $\theta_2$ or
injectivity of multiplication by monic polynomials on its source.
\end{saresult}
\begin{proof}
Write $A=\Omega(R,S)$. By B.064 every coefficient of $A$ lies in $C$,
and injectivity of $\theta_2$ makes $R\to C$ epic.
The multiplication map $\mu:C\otimes_R A\to C[X]$ is always surjective,
because $R[X]\subset A$. Suppose that
$z=\sum_{i=1}^m c_i\otimes F_i$ lies in its kernel.
For each $c_i$, B.064 supplies a monic $P_{d_i}\in R[X]$ such that
$c_iP_{d_i}\in A$. Their product $H=\prod_iP_{d_i}$ is monic and
$c_iH\in A$ for every $i$. The tensor
\[
                         w=\sum_i c_iH\otimes F_i
\]
has polynomial image $H(X)\sum_i c_iF_i(Y)=0$. Injectivity gives $w=0$.
Apply extraction of the coefficient of $X^{\deg H}$ in the first
factor, as an $R$-linear map $A\to C$. The result is $z=0$, proving
the first isomorphism in (B79.2).

Reduction of coefficient vectors modulo $R$ gives an exact sequence
\[
                0\longrightarrow R[X]\longrightarrow A
                    \longrightarrow N(R,Q)\longrightarrow0.  \tag{B79.3}
\]
Surjectivity on the right follows by lifting the finitely many
coefficients of a vector in (B79.1): its defining equations say
precisely that the lifted polynomial has values in $R$.
After tensoring, the composite
$C[X]=C\otimes_RR[X]\to C\otimes_RA\to C[X]$ is the identity.
Thus the first arrow is split injective and its cokernel is
$C\otimes_RN(R,Q)$. The second arrow is an isomorphism exactly when
this cokernel vanishes, proving the equivalence claimed.

Here is the explicit obstruction to sufficiency and to monic
saturation. Let $k=\mathbb F_q\subset l$ with $l$ infinite, let
$V=l[[t]]$, $K=l((t))$, $R=k+tV$, and $S=K$. Put $P=X^q-X$ and choose
$\beta\in l\setminus k$. Formula (B77.4), equivalently (B78.7), gives
the direct $R$-module decomposition
\[
 A=\bigoplus_{i<q}RX^i\ \oplus\!
      \bigoplus_{\substack{i<q\\j\ge1}}t^{1-j}V X^iP^j.
                                                        \tag{B79.4}
\]
The effective coefficient ring is $K$: for every $n\ge0$ and
$v\in V$, the polynomial $t^{-n}vP^{n+1}$ has values in $tV\subset R$.
Consequently $C=K$ is a flat epic localization, and
$K\otimes_RA=K[X]$.
Nevertheless
\[
                       z=\beta P\otimes P-P\otimes\beta P
                                                        \tag{B79.5}
\]
is a nonzero tensor with zero polynomial image. Its projection using
the coefficient of $P$ in (B79.4), followed by reduction modulo $tV$,
is $\beta\otimes1-1\otimes\beta$ in $l\otimes_k l$, which is nonzero.
For example choose a $k$-linear functional on $l$ that vanishes on
$1$ and takes value one on $\beta$.
But $P^2/t\in A$, since $P(R)\subset tV$, and $t,t\beta\in R$.
The balancing relations in the tensor product give
\[
 \beta P^2\otimes P
       =(P^2/t)\otimes(t\beta P)
       =P^2\otimes\beta P.
                                                        \tag{B79.6}
\]
Therefore $P(X)z=0$. This is a concrete nonzero element killed by a
monic polynomial in the first variable, despite flat epic coefficient
base change and both conditions (B79.2). It prevents that necessary
condition from closing the unrestricted structural question.
\end{proof}
